\documentclass[11pt,a4paper]{article}
\usepackage[T1]{fontenc}
\usepackage{XCharter}
\usepackage[margin=24mm]{geometry}
\usepackage{microtype,amsmath,amssymb,booktabs,tabularx,longtable,array}
\usepackage{xcolor}
\usepackage{xurl}
\definecolor{accent}{RGB}{21,102,121}
\usepackage[colorlinks=true,linkcolor=accent,urlcolor=accent,citecolor=accent]{hyperref}
\setlength{\emergencystretch}{2em}
\setlength{\parskip}{3pt}
\renewcommand{\arraystretch}{1.15}
\newcommand{\pair}[1]{\texttt{#1}}
\newcommand{\ev}[1]{{\small\textsf{Evidence:} #1}}
\newcommand{\sel}{\(\bullet\)}
\title{Recursive Science\\\large What repeats and reinjected papers added to a 10-by-10 Pathfinder experiment}
\author{Pathfinder experiment record}
\date{28 September 2026}

\begin{document}
\maketitle
\begin{abstract}
Can a research pipeline use its own completed papers as productive inputs?
We examine a Pathfinder experiment beginning with ten papers returned for
``mechanism design'' and ten for ``agentic cooperation''. After investigating
14 selected pairs and repeating those investigations, we compared another
14-pair repeat with 14 pairings of original papers and previously accepted
manuscripts. All 28 pilot pipelines finished. The repeat arm produced seven
internally accepted papers and reinjection three, but claim-level comparison
changes the interpretation: several repeat papers reconstruct earlier work,
including a payment repair already developed in a separate follow-up.
Reinjection produced a concrete correction and graph extension of a
peer-ranking seed, and a conditional-risk counterexample crossing an
obedience boundary. Both arms developed related auction information limits.
The scientific objective is extension of inquiry: supported ideas developed
through generated intermediate results and not observed in the historical
or matched repeat record. It is not superiority over repetition, nor a claim
that repetition could never discover those ideas. The useful recursive
unit is a claim with assumptions and ancestry, not an accepted PDF alone.
\end{abstract}

\section{The question: repetition is not recursion}

Pathfinder searches for possible connexions between two collections of papers.
An abstract-level scan proposes a question for each pair. Selected pairs pass
to two researchers, who inspect sources, exchange findings and objections in
a ledger, and produce a consolidated note. A verifier may request further
work, stop the thread, or recommend a paper. An editor produces a readable
account; a recommended paper then passes through author and reviewer stages
\cite{engine}.

A repeat gives a fresh investigation the same original inputs. Reinjection
replaces an input by a completed, internally accepted paper from an earlier
investigation. Here reinjection occurs at the very end of the pipeline,
after paper writing and review, not immediately after editing a research note.
It exposes a later team to a compressed argument, including its assumptions,
proofs, references and unresolved questions.

Different answers on a second run are not evidence for recursion: fresh
research on unchanged inputs also follows different paths. Waiting until
repeats appear exhausted would not solve this problem cleanly. There is no
known stopping rule that certifies exhaustion, and waiting changes the
resources spent before recursion is tried. We therefore used a contemporaneous
repeat arm as a control, without claiming saturation \cite{protocol}.

\section{The original 10-by-10 space}

The saved corpora contain ten Q papers and ten P papers, identified in
Appendix~\ref{app:corpus}. These are query-returned collections, not pure
disciplinary samples: Q includes robotics and astronomical instrumentation
as well as economic mechanisms. The experiment studies this particular
retrieved space, not a random sample of science.

The scanner scored feasibility and scientific gain from titles and abstracts,
each on a 0--100 scale, and ranked their product. The archive contains 100 scan
records, of which 99 were scored successfully. A product threshold of 1500
selected 13 pairs; the human-retained Q1P2, scored 600, made the research
shortlist 14. The remaining 86 cells were not investigated in these campaigns.
Q1P2 later produced an accepted selective-perception paper. This is evidence
of a useful exception to the scan, not an estimate of its false-negative rate.

\begin{table}[ht]
\centering\small
\begin{tabular}{l*{10}{c}}
\toprule
 & P1&P2&P3&P4&P5&P6&P7&P8&P9&P10\\
\midrule
Q1&\sel&H&&&&\sel&&&&\\
Q2&&&&&&&&&&\\
Q3&&&\sel&&&\sel&&&\sel&\sel\\
Q4&&&\sel&&&\sel&&&\sel&\sel\\
Q5&&&&&&&&&&\\
Q6&&&&&&&&&&\\
Q7&\sel&&&&&&\sel&&&\\
Q8&&&&&&&&\sel&&\\
Q9&&&&&&&&&&\\
Q10&&&&&&&&&&\\
\bottomrule
\end{tabular}
\caption{The investigated subset of the original space. Dots mark the
13 threshold selections; H marks the retained human exception. These are
selection marks, not successful-result marks.}
\end{table}

The original campaign used mixed models and subsequent interventions; its
receipts include Claude and GPT stages. A September 23 fresh-context rerun
used GPT-6-sol on the same 14 pairs. Their outcomes already showed substantial
variation: the original produced five accepted papers and the rerun three.
Q4P6 moved from a conditional peer-calibration construction to a PAUSE on
strategic shaping; Q8P8 moved from a proposed compliant-grasp study to a
conditional wrist-observation obstruction. A changed verdict need not refute
an earlier argument: it may reflect a different argument or evidential bar.

\section{The repeat--reinjection pilot}

Five final papers were frozen as seeds, selected for different argument
families rather than solely for recency. Their full text and compiled
bibliographies were preserved, with artifact hashes and parentage in the
manifest \cite{pilot}.

\begin{table}[ht]
\centering\small
\begin{tabularx}{\linewidth}{llX}
\toprule
Seed&Parent&Argument carried forward\\
\midrule
S1&O/Q4P6&Shared-comparison calibration and contribution ranking.\\
S2&O/Q1P2&Room-level calibration transferred to a selective wall gate.\\
S3&R/Q3P10&Auction surplus with explicit computation costs.\\
S4&R/Q7P1&True versus expected empirical CVaR, with conditional economic bounds.\\
S5&R/Q8P8&Limits of wrist-only information for an occluded held object.\\
\bottomrule
\end{tabularx}
\caption{O denotes the original campaign; R the September 23 rerun.
The carbon-mechanism paper was deferred because its source assumptions needed
audit. S4 was carried forward for its risk distinction, not as certification
of the source mechanism.}
\end{table}

The control C repeated the original 14 pairs. The reinjection arm I used
14 distinct Q--seed pairs, excluded each seed's parent Q, and matched Q
frequencies exactly: Q1 appeared three times, Q3 four, Q4 four, Q7 twice,
and Q8 once. Scheduling blocks alternated which arm ran first. No new scan
selected treatment pairs after their outcomes were known. In I, local P1--P5
mean S1--S5, not original P1--P5.

Both arms used GPT-6-sol at high reasoning effort, the same frozen engine
and prompts, two peer seats, up to four research rounds, and up to three
paper rounds. They had separate USD 200 API-equivalent admission caps.
Equal caps did not require equal realised expenditure. Each investigation
could stop at PAUSE; only DRAFT proceeded to a paper. Only one recursive
generation was authorised. Pilot outputs have not been reinjected again.

There were operational interruptions, including authentication and transport
failures. Resumption retained stage identities and completed artifacts rather
than silently treating retries as new scientific replicates. The final
runner and supervisor both report completion; all 28 edits and ten papers
report successful builds. All 720 artifacts protected at the last resumption
were unchanged, and frozen-input preflight passed. These integrity checks do
not establish mathematical validity or complete reference verification.

\begin{table}[ht]
\centering\small
\begin{tabular}{lrrrr}
\toprule
 & O & R & C & I\\
\midrule
Investigated pairs &14&14&14&14\\
Internally accepted papers &5&3&7&3\\
Other terminal research outcomes &9&11&7&11\\
\bottomrule
\end{tabular}
\caption{Descriptive pipeline outcomes. O includes one PAUSE-ON-ITERATE
among its nine other outcomes. O and R are historical context, not controlled
model comparisons. C and I form the contemporaneous pilot.}
\end{table}

Recorded arm costs were USD 52.03 for C and USD 36.52 for I, from 119 and
105 receipts respectively. C has two error receipts. Usage is missing from
one of them. I has no error receipts. These are recorded API-equivalent workloads under subscription
access, not additional invoices or guaranteed complete metering. They exclude
the historical cost of seeds \cite{readout}. Supervision is separate: 115 audit logs contain
114 completed-turn usage records, totalling 24,120,478 input tokens, of which
18,601,984 were cached, and 244,038 output tokens. Failed calls and separate
probes may lack usage. The pilot took 49.24 hours from staging to completion,
including interruption; summed model-call duration is not elapsed wall time.

\section{How claims were compared}
\label{sec:method}

The comparison unit is a statement with assumptions and an argument, not a
paper, title, verdict or topic. We distinguish inherited results, extensions,
corrections, new connexions, and proposals still lacking decisive evidence.
``Supported'' here means that the recorded derivation or finite computation
supports the stated conditional claim; it does not mean externally reviewed,
experimentally demonstrated, or new in the literature.

The main historical baseline is the union of O and R, including PAUSE ledgers.
Close comparisons also use the already-existing Q7 cache-run, September 13
paper follow-ups and September 14 source audit. This is an explicit extension
of the initial two-campaign inventory: otherwise the pilot would receive
credit for a payment repair already present in the repository. Other model
benchmark campaigns are not exhaustively assessed here. Thus novelty means
``not found in this specified record'', not ``first found anywhere''.

We inspected terminal claims and the relevant proofs, compared ledger
entries and earlier supporting notes, and searched 186 historical authored
and ledger files for the distinctive constructions. Search matches were
checked for substance: a source line numbered 396 is not a 396-board theorem.
Absence of a phrase was used to guide reading, not as proof of novelty.
The all-board witness was rerun successfully; the common-flip endpoint
calculation was independently checked numerically. This is not a fresh
line-by-line proof audit of every manuscript.

The planned label-hidden support assessment was not performed: this reviewer
had already seen arm labels. The classifications below are therefore a
transparent, unblinded post-hoc assessment. The same distinction between
proved conditional results and untested applications is applied to both arms.
Appendix~\ref{app:inventory} records a disposition for every pilot pair,
including PAUSE. Granularity affects counts, so we report result families
and their differences rather than manufacture a total ``discovery score''.

\section{What survives deduplication}

\subsection{Peer rankings: a real recursive extension, not the first warning about bias}

S1's positive construction uses reports $Y_{er}=X_eZ_{er}$: observers share
one realised comparison $X_e$, with conditionally independent symmetric
errors of positive means $a_r$. Shared-item agreements satisfy
$q_{rs}=a_ra_s$. Calibrating observer reliabilities then permits recovery
of Bradley--Terry score differences under coverage assumptions.

The historical record already warned that well-fitted comparison scores need
not measure true contribution. R/Q3P6, for example, lets a nuisance cue shift
the latent score while preserving a Bradley--Terry form. The September 13
S1 exposition also explicitly warns that passing observable checks does not
establish conditional independence. Those warnings are inherited.

I/Q1P1 adds a more specific construction. An independent common flip $B_e$,
with constant mean $b>0$, gives $Y_{er}=X_eB_eZ_{er}$. Agreements still obey
$q_{rs}=a_ra_s$, while calibrated edge means become
$b\tanh((c_i-c_j)/2)$. A four-vertex path can therefore have the same full
report law in two worlds with opposite endpoint rankings. At $b=1/2$,
the true path differences $(4,-3/2,-3/2)$ sum to 1, but their uncorrected
transforms sum to approximately $-0.26458$.

This supplies a counterexample to the seed's broad sentence that correlated
errors destroy factorisation; it does not refute the seed's theorem under
conditional independence. The new graph results are also substantive:
diameter at most two preserves ordinal information within the specified
common-flip model; suitable triangles or four-cycles identify its missing
scale, while a recorded six-cycle admits ambiguity. These results and a
conditional finite-sample interface are not in the inspected historical
baseline or C. The general bias warning is old; the explicit same-law
construction and graph boundary are the incremental result.

This paper uses S1 deeply. Q1 contributes the distinction between ordinal
and numerical predictions and their downstream cost, rather than a new
algorithmic performance theorem. Its relevance is principally a correction
and extension of the seed, not equal technical contributions from both inputs.

\ev{O/Q4P6 ledger 28, 32, 37, 41; R/Q3P6 ledger 7, 9;
September 13 Q4P6 exposition, limitations;
I/Q1P1 ledger 4, 8, 11, 16, 26, 27 and final paper.}

\subsection{CVaR and obedience: a new conditional connexion}

S4 already establishes that the expectation of empirical CVaR can differ
from true CVaR and reverse an economically relevant comparison. I/Q4P4
does not discover that bias again. Its increment is to connect the bias to
Q4's state-independent obedience-cycle test, rather than the seed's resource
tax game. The researchers first identify that the paper called Q inside
the seed is not the Q assigned to this new pair.

For each finite batch size $m$, a risky loss is 0 or $2h$ with equal
probability. Its expected empirical upper-half CVaR is at most
$2h(1-2^{-m})$. Put a safe loss $c_m=2h-h2^{-m}$ between that value and
true CVaR $2h$. In two signals with exchanged action losses, the surrogate
supports the prescribed risky actions with zero rewards. True obedience
would require both $r(a)-r(b)\geq h2^{-m}$ and its reverse with the same
positive margin, which is impossible. The same policy thus crosses an
implementability boundary, not merely an objective-value boundary.

This is a supported new connexion relative to the inspected baseline and
C, where the risk analysis stays with resource allocation or query cost.
It is deliberately narrow: the instance changes with $m$; risk values are
conditional reduced-form utilities; the impossibility concerns
state-independent rewards, not unrestricted state-dependent rewards or
CVaR of an entire plan. Cycle-margin and common-reward certificates are
useful extensions of standard bounded-error and difference-constraint ideas.

\ev{R/Q7P1 paper, fixed-batch participation example and value bounds;
I/Q4P4 ledger 1, 6, 9--16 and final paper.}

\subsection{Auction baselines: useful sharpening, substantial inheritance}

I/Q1P3 inherits the surplus ceiling 7.5 from S3. Even earlier, O/Q3P3
computed trade-count ceilings, omission/commission losses, and stylised
zero-intelligence benchmarks. A claim that reinjection discovered the need
for a cheap auction baseline would therefore be wrong.

The new paper supplies a specific conditional no-generated-token fixed-price
rule and a price-error counterexample. If its five eligible pairs execute,
the rule reaches 7.5. An arbitrarily small upward advised-price error can
admit a cost-2 seller who displaces a cost-0.75 seller, losing 1.25 surplus.
The latter distinguishes scalar price accuracy from allocation quality.
Neither construction establishes behaviour under the actual random scheduler.
This is a modest, useful extension of an existing family, not a wholly new
auction result or measured benefit of advice.

C/Q3P3's accepted paper also belongs to that family. Its five/six/seven-trade
accounting sharpens an older warning that more individually profitable
crossings need not increase welfare. I/Q3P2, although PAUSE, adds a
rounding-aware lower bound: at least about 44.9\% of a reported model's
aggregate deficit is attributable to trade count. That numerical
interpretation is not found in the historical baseline, but is Q-only:
the calibration seed does not supply the bound. Its proposed market
calibration lacks run-level data. These should not become three independent
discoveries of the same welfare-accounting principle.

\ev{O/Q3P3 ledger 3--9, 12; R/Q3P10 paper;
C/Q3P3 final paper; I/Q1P3 ledger 10, 15, 19, 24;
I/Q3P2 ledger 14, 16, 18, 22, 24.}

\subsection{Auction information limits: a shared family, with a meaningful difference}

C/Q3P6 constructs two worlds with the same acting-trader information and
the same value multiset, but opposite welfare rankings of crossing and
improving a bid. Its final-step continuation is stipulated. Under an equal
prior, no local judge can exceed one-half accuracy on the realised better
action; the illustrated value of revealing a hidden buyer value is 0.25.
This is an explicit auction instantiation beyond O/R's comparator-bias
proposals, not evidence about the frequency of such states.

I/Q3P5 imports S5's matched-observation idea into the auction and reaches a
related result. Its ledger rejects an initial rent-only example and a
symmetry problem, then swaps a value between an active and an already-traded
agent and averages over future random invitations. Thus it extends the
construction beyond conditioning on a particular next actor. It still needs
a stipulated continuation and illustrative prior.

The pair of outputs belongs to one auction-information family, with different
conditional examples. The repeat was ACCEPTED and reinjection PAUSE, despite
both offering mathematical examples with unresolved empirical relevance.
Counting only papers would conceal the shared finding and this threshold
variation. Reinjection adds a refinement here, not exclusive ownership of
the information-limit idea.

\ev{O/R Q3P6 ledgers; S5 final paper; C/Q3P6 ledger 19, 20, 22, 28;
I/Q3P5 ledger 4, 15, 18--21.}

\subsection{Contracts: the clearest repeat-only constructive addition}

O/Q4P9 focused on representation and enforcement confounds. R/Q4P9 then
proposed path-based contract sufficiency and distinguished feasibility from
obedience. C/Q4P9 goes further: its enumeration supplies a reciprocal
five-clause tile contract and shared shortest route for every one of the
396 asymmetric boards, giving scores $(85,55)$ or $(75,65)$ instead of
the outside scores $(70,0)$ and achieving joint score 140.

We reran the archived enumeration: 314 boards need one Red coverage and
82 need two, reproducing the score counts. This verifies the finite
construction under the program's model, not source-code equivalence with
every simulator rule. It is a supported extension beyond the historical
proposal and was not found in I. It proves an available mutually improving
outcome, not voluntary bargaining or sequential equilibrium. Q supplies
the feasibility/obedience distinction; the computation is predominantly
about the contract game's rules.

\ev{R/Q4P9 ledger 14, 17, 19; C/Q4P9 ledger 20--26,
final paper and \path{paper/asymmetric_oracle.py}.}

\subsection{Payment repairs and reward shaping: accepted does not mean new}

C/Q7P1's explicit others-only transfer correction looks new against a short
summary of O and R. It is not new against the existing follow-up record.
The September 13 paper already has the same correction: its coefficient
$\gamma(b_N-1)/(N-1)$ equals $\gamma N/(N-1)^3$, the coefficient in C.
The September 14 source audit also supplies direct allocation implementation,
participation, non-unique prices and risk-surrogate error control. C's
rediscovery is useful confirmation and exposition, not another independent
repair. C/Q7P7 combines that repair with a previously developed carbon charge
and forecast-error argument, under restricted optional loads. Its synthesis
must not count the same repair again. I/Q7P2's source contradictions also
repeat earlier findings; its calibrated dual-gap proposal remains conditional.

C/Q1P6 is similarly mostly inherited. O/Q1P6 already proves fixed-population
telescoping, entry/exit accounting and the loss of ranking information when
softmax potentials are pooled. R/Q1P6 supplies critic-offset cancellation and
the finite unregularised Bradley--Terry estimate counterexample. Recombining
these in an accepted paper is not a new finding that comparisons do not
automatically improve learning.

\ev{September 13 Q7P1 paper, repair equation;
September 14 Q7 source audit, static repair;
O/Q1P6 ledger 24, 27, 30, 33; R/Q1P6 ledger 4, 9, 18;
C/Q7P1, C/Q7P7 and C/Q1P6 final papers.}

\section{Focused mathematical audit}

After the claim comparison, we rederived the principal ranking and risk
arguments and checked their numerical witnesses with
\path{notes/check_recursive_claims.py}. This is a separate calculation by
the same supervising agent, not blinded external review, literature
verification or an empirical test. Archived campaign outputs were left
unchanged.

\paragraph{Ranking witnesses and proof boundaries.}
The path example gives intended endpoint difference $1$ and transformed
difference $-0.264580326901$. For the six-cycle, closure reduces to
$S_1b^4+S_3b^2+S_5=0$, where $S_k$ is the degree-$k$ elementary symmetric
sum of the six calibrated means. Its admissible roots are
$b=0.738908509686$ and $b=0.853764730071$; both close the cycle and reverse
the order of vertices 3 and 6, numbering from the first listed edge.
These checks support the claimed nonidentification, not just a poor fit.
The diameter-two argument follows from the sign identity for an odd,
strictly increasing transform; the necessity statement remains restricted
to trees. The triangle and four-cycle identities require nondegeneracy.
The finite-sample bound follows by propagating log-agreement error,
bounding the logit derivative and applying the centred incidence-matrix
pseudoinverse. It still assumes a bound on bias relative to the clean
target; more labels cannot certify that assumption.

\paragraph{Risk reversal and reward certificates.}
For loss $0$ or $2$ with equal probability and upper-half tail risk, the
empirical value with $k$ high observations among $m$ is
$\min(2,4k/m)$. Exact binomial summation for $m=1,\ldots,64$ confirms the
paper's inequalities, including odd batch sizes. The general argument is
the all-zero batch event of probability $2^{-m}$, so the proof is not
limited to those checked sizes. The conditional-versus-ex-ante example
also reproduces the reported ex-ante value $3.194>3$.
The cycle error bound uses two utility errors per signal; the rectangular
robust-reward criterion is the usual feasibility condition for difference
constraints. No defect was found in these arguments under their stated
assumptions. The finite-batch instance changes with $m$, and neither
report-stage incentives nor realised-batch guarantees follow.

\paragraph{Auction arithmetic passes, but the waiting policy needs precision.}
C/Q3P6's conditional information value is $0.25$. I/Q3P5's displayed
random-invitation regret expression evaluates to $0.112198829651$.
However, the latter's supporting construction allows the buyer to cross
again if invited while also arguing that selection of the cheap seller
within six draws guarantees surplus $2$. Early crossing by the buyer can
remove that opportunity, so the lower-bound argument needs a more explicit
policy than merely avoiding loss-making trades.

A sufficient clarification is to let the informed comparator in state A
keep its bid of $2$ standing for all six draws, with other agents holding
their quotes except for the cheap seller's profitable crossing. Then the
claimed $2[1-(7/8)^6]$ lower bound holds. The state-B upper bound for a
noncrossing first action still follows because a subsequent profitable
trade requires the buyer's reinvitation. Together with the state-A upper
bound $0.75$ after immediate crossing, these yield the recorded regret
bound under this explicit continuation. Thus the information-limit
construction is repairable, but its written policy description should be
tightened before reuse. We have checked the arithmetic and this logical
qualification, not implemented a simulator-equivalent replay or measured
the frequency of the constructed history.

\section{Conclusions for recursive science}

\paragraph{Recursion can change the object of inquiry.}
The ranking seed becomes an object to challenge, not merely an authority to
reuse. The risk seed supplies an ingredient for a different implementability
question. These are concrete reasons to continue studying reinjection. They
do not show that either result required a generated input or could not have
been found by an unobserved repeat.

\paragraph{The objective is extension, not replacement.}
The repeat arm finds a new finite contract witness and its own information
construction. Reinjection contributes distinctive extensions but fewer
accepted papers. These counts do not determine whether inquiry was extended.
The relevant evidence is a supported new claim, its dependence on a generated
intermediate result, and its difference from the historical and control
findings. The ranking and risk examples meet that descriptive criterion in
the inspected record. A finite repeat sample cannot establish that repetition
would never produce them. Shared seeds, selected pairs, heterogeneous research
paths and incomplete blinding also rule out a universal causal conclusion.

\paragraph{The historical ledger matters as much as the final paper.}
PAUSE can contain a sound bound or counterexample whose application is
unsupported. Earlier PAUSE work prevents false credit for later acceptance.
Separate follow-ups matter too: a paper-only or two-directory inventory would
misclassify the payment repair. Ancestry should track claims, assumptions,
corrections and unresolved dependencies, not just parent filenames.

\paragraph{Acceptance is a staging decision, not an epistemic guarantee.}
The seed's overbroad correlation statement survived internal review. The
recursive paper usefully challenges it. Conversely, repeating an accepted
argument can amplify its authority without adding evidence. All pilot
editors also report an arXiv HTTP 406 lookup failure; local generated seeds
have missing-public-identifier warnings. Successful builds do not resolve
those checks. Reference verification and independent mathematical review
remain separate work.

\paragraph{A focused continuation is preferable to waiting for saturation.}
The focused audit above supports the common-flip and CVaR-cycle results
under their restricted models and identifies a continuation-policy
qualification in the recursive auction example. External mathematical
review remains open. The September 28 source audit verified the selected
seeds' public references and checked their attribution boundaries, including
the source signal-cycle lemma. It is not an exhaustive prior-art search.
Versioned seed copies apply the archived reviewers' minor corrections;
the historical outputs are unchanged. A separate auction corrigendum supplies
an explicit sufficient waiting policy.

The next protocol selects two generated seeds and two original-input controls,
with the same incremental question within each block. It is a purposive
second-generation extension test, not a contest in acceptance rates. It does
not feed every accepted PDF into an ever-growing pool. Its allocation,
admission limits and subsequent run status are recorded separately in
\path{plans/2026-09-28-1322-recursive-extension.md} and
\path{experiments/2026-09-28-recursive-extension/}; no outcome from that
follow-up is included in the first-pilot conclusions above.

\ev{\path{plans/2026-09-28-source-audit.md};
\path{notes/recursive-auction-corrigendum.tex}.}

Daily query intersections, absolute-interest filtering and broader corpus
growth are possible later selection strategies, not interventions tested by
this pilot. The present result is narrower: a fixed small seed space supports
both rediscovery and substantive transformation, and a claim-level memory is
needed to tell them apart.

\appendix
\section{Complete pilot disposition}
\label{app:inventory}

The tables give the residual after comparison, not a binary scientific-quality
score. A modest extension can be mathematically sound yet insufficient as a
paper claim. ``New'' is relative to the inspected record described in
Section~\ref{sec:method}. The primary evidence for each row is its terminal
paper or edited note and ledger; the named O/R antecedents provide the
comparison. A means internally accepted paper; P means PAUSE.

\small
\begin{longtable}{p{17mm}p{7mm}p{112mm}}
\toprule
Repeat&End&Claim-level disposition\\\midrule
\endfirsthead
\toprule Repeat&End&Claim-level disposition\\\midrule\endhead
Q4P10&A&Extension: split-prize collision calculation and a bounded priority
realisation. The ordered-protocol theorem is established prior work, as the
paper acknowledges. O/R pursued paid peer scoring instead. No survival gain
was measured.\\
Q3P9&P&Mostly repeated proposal: immediate settlement is not missing
commitment; proposal order and outside options remain confounded. O/R Q3P9
already separate commitment from wording. A role-by-order test is unrun.\\
Q4P6&P&Conditional extension: bounded payments must distinguish every mixture
of profitable deviations, not just each deviation separately. The finite
monitor/minimax test and shirking benchmark sharpen O/R's shaping boundary;
no strategic camera or success law was measured.\\
Q3P10&P&Repeated accounting and untested tariff proposal. The 7.5 ceiling,
first-call and scheduler cautions already occur in O Q3P3 and R Q3P10.\\
Q7P1&A&Rediscovery against the expanded baseline: the same others-only
repair and static risk comparison are in the September 13--14 follow-ups.
Different presentation is not a new mechanism.\\
Q4P3&P&Modest extension: two-way separation of peer informativeness and defect
truth, plus an F1 filtering inequality. O/R already warn that agreement and
copied evidence are insufficient. No review intervention was tested.\\
Q1P1&P&Extension of R/Q1P1's rare-tail lower bound from value estimation to
an optimisation example with opposite optima. Threshold lifting is known;
no new prediction-sensitive query saving is established.\\
Q3P3&A&Refinement of old welfare-accounting family, not a new discovery of
7.5 or of inefficient extra crossings. O/Q3P3 already records ceilings and
omission/commission. The review intervention remains untested.\\
Q4P9&A&Supported finite extension: all-board reciprocal contract witness
and delivered-coverage identity, beyond R/Q4P9's path-audit proposal.
Predominantly P-side, not an incentive theorem.\\
Q7P7&A&Recombination: earlier payment repair plus O/R's carbon charge,
restricted allocation and forecast-error arguments. Conditional synthesis
for optional loads; not an additional independent repair.\\
Q8P8&P&Formal extension: identical Gaussian observations can hide opposite
causal directions; a scalar pose-spread statistic can hide free yaw. O/R
already question rigidity as a proxy for causal attachment. No hand failure
or control improvement was measured.\\
Q3P6&A&New explicit auction-information example, sharing a family with
I/Q3P5. Fixed value multiset, stipulated continuation, no frequency estimate.\\
Q1P6&A&Largely reconstruction: O/Q1P6 has pooling and entry/exit accounting;
R/Q1P6 has critic cancellation and estimator separation.\\
Q1P2&P&Old first-gate margin plus an elementary unit-audit lower bound
$n(1-2\delta)$. O/R already isolate missing observations and fallback
conditions. The bound is not a physical LiDAR-cost certificate.\\
\bottomrule
\end{longtable}

\begin{longtable}{p{17mm}p{7mm}p{112mm}}
\toprule
Reinject&End&Claim-level disposition\\\midrule
\endfirsthead
\toprule Reinject&End&Claim-level disposition\\\midrule\endhead
Q4P2&P&Elementary extension of S2: calibrate a maximum over a fixed fitter
menu to resist selection. Generic simultaneous coverage and the two-action
obedience inequality are not a new measured wall-gate result.\\
Q3P1&P&Inherited shaping invariance; hypothetical payoff-equivalent feedback.
Public-tape ambiguity and a common-anchor example sharpen cautions, but do
not establish a feedback effect. Relates to O/R Q3P3 and Q3P6.\\
Q4P3&P&Modest tariff extension: own-token charges can miss counterparty
costs in a fixed-match subgame. S3 already identifies avoided calls.
The corrective reward needs metering, funding and enforcement.\\
Q3P2&P&Supported Q-only numerical extension: the 44.9\% deficit bound.
Welfare ceilings are inherited. The S2-style calibrated market guarantee
has no run-level data and is not a new calibration theorem.\\
Q7P3&P&Conditional accounting extension: paid learning can exhaust the
participation surplus; a geometric sample schedule gives an iteration cap.
S3 supplies cost accounting, but Q7 does not posit paid LLM samples.\\
Q4P4&A&New connexion: finite-batch conditional CVaR reverses an
obedience-cycle test. Inherits S4's surrogate bias, adds the specific
state-independent implementability boundary and error certificates.\\
Q1P3&A&Modest extension: completed fixed-price comparator and a
price-error/rationing example. S3 and O/Q3P3 already supply the welfare
ceiling and cheap-baseline motivation.\\
Q3P4&P&Conditional instantiation of S4's risk reversal in an auction
cross/wait choice. The stipulated continuation and analyst-only fill
probability do not explain observed trader behaviour.\\
Q4P5&P&Extensions under added models: observation-testing bound,
measurement-based incentive obstruction, and a coarse-signal oversight
counterexample/repair. S5 supplies the information issue. No matched-history
handover data or strategic donor is established.\\
Q7P2&P&Source corrections are inherited from R/Q7P7 and the September 14
audit. S2-inspired dual-gap calibration is a conditional proposal with no
trusted episode labels or nonvacuity demonstration.\\
Q8P4&P&Elementary S4 risk example and a zero-fall sample calculation
requiring 598 independent episodes under its stated target. No closed-loop
hardware losses establish a controller result.\\
Q3P5&P&Supported auction-information construction; shares C/Q3P6's family
but integrates over random invitations. Refinement, not an exclusive
reinjection discovery or observed explanation of market losses.\\
Q1P5&P&Mostly S5's inherited obstruction, extended with a standard
value-of-information identity and marginal set coverage. Practical sensing
benefit remains unmeasured.\\
Q1P1&A&Correction and substantive extension: shared flips can preserve
agreement factorisation while reversing rankings; graph structure separates
ordinal recovery and scale identification. Seed-driven, with Q1 framing
the prediction interface.\\
\bottomrule
\end{longtable}
\normalsize

\section{Corpus identity and evidence map}
\label{app:corpus}

The abbreviated corpus labels below preserve the original ordering. Full
titles, abstracts and identifiers are in the root \path{Q.jsonl} and
\path{P.jsonl}. The query-date prose in the engine note and the earliest
saved scan receipts do not agree; corpus files and frozen hashes, rather
than that date, identify the experiment here.

\begin{longtable}{lp{108mm}}
\toprule ID&Paper (arXiv identifier)\\\midrule\endfirsthead
\toprule ID&Paper (arXiv identifier)\\\midrule\endhead
Q1&Learning-Augmented Algorithms (2609.04787)\\
Q2&Approximately Efficient Multidimensional Bilateral Trade (2609.02872)\\
Q3&Competitive Market Behavior of LLMs (2609.02580)\\
Q4&Mechanism Design for Alignment and Control (2609.01595)\\
Q5&Sensitivity Hot Spot Penalization (2608.30499)\\
Q6&Facility Location Games Under a Prelocated Facility (2608.30292)\\
Q7&Mechanism Design with Stochastic Dynamic Stability (2608.29130)\\
Q8&Aero Hand Open (2608.28578)\\
Q9&Optomechanical inertial reference for atom interferometry (2608.28418)\\
Q10&Canadian Galactic Emission Mapper (2608.28159)\\
P1&Distributed risk-averse optimization via CVaR (2609.04460)\\
P2&Concept-Guided Exploration (2608.23650)\\
P3&Adversarial Review (2608.18167)\\
P4&Uplink-Completion-Triggered Edge-GPU Inference (2608.08330)\\
P5&Multi-Agent Reinforcement Learning for Base Station Placement (2607.28002)\\
P6&MARS-RA (2607.27967)\\
P7&Emission-Forecasting-Based Spatial-Temporal Carbon Response (2607.26560)\\
P8&One Hand Watches The Other (2607.22119)\\
P9&Commitment To Cooperation With Self-Negotiated Contracts (2607.22750)\\
P10&The Energy Society (2607.14865)\\
\bottomrule
\end{longtable}

Evidence paths are relative to the repository root. O/pair means
\path{threads/<pair>/}; R/pair means
\path{experiments/2026-09-23-gpt-6-sol-rerun-01/threads/<pair>/}.
C and I mean respectively \path{repeat/threads/<pair>/} and
\path{reinjection/threads/<pair>/} beneath
\path{experiments/2026-09-25-repeat-injection/}.
Within each thread, \path{ledger.jsonl} records numbered entries;
\path{paper/paper.tex} is the final paper when present, and
\path{edited/note.tex} is the readable account. Verdicts do not replace
either source of evidence.

The expanded historical baseline additionally includes
\path{experiments/cache-run/threads/Q7P1/},
\path{experiments/2026-09-13-paper-followups/threads/}, and
\path{experiments/2026-09-14-Q7-source-audit/Q7-repair.tex}.
The latter audit and the September 13 Q7P1 paper are essential to the
payment-repair deduplication. Their inclusion changes the novelty judgment,
not any frozen campaign input or outcome.

The pilot manifest preserves exact seed ancestry, hashes, corpus and prompt
hashes, configuration and the 28-entry schedule. Its \path{outcomes.json},
arm \path{receipts.jsonl} files and final supervisor audit support the
operational totals. These archived records remain unchanged by this review.

\begin{thebibliography}{9}
\bibitem{engine} Pathfinder. \emph{The Pathfinder engine}.
Repository note, \path{notes/pathfinder-engine.tex}.
\bibitem{protocol} Pathfinder. \emph{Repeat versus reinjection: claim
baseline and pilot protocol}. 25 September 2026.
\path{plans/2026-09-25-1623-repeat-versus-reinjection.md}.
\bibitem{pilot} Pathfinder. \emph{Repeat--reinjection pilot archive}.
25--27 September 2026. \path{experiments/2026-09-25-repeat-injection/}.
See Appendix~\ref{app:corpus} for the thread-level evidence map.
\bibitem{readout} Pathfinder. \emph{Repeat versus reinjection: first readout}.
27 September 2026. \path{plans/2026-09-27-repeat-reinjection-readout.md}.
Its initial candidates are superseded by the historical comparison here.
\end{thebibliography}
\end{document}
